\subsection{The joint cubic--quartic cyclic cover}

\begin{statement}{D23.F21}{Proved}{The degree-twelve cyclic cover is blind at every two-adic and three-adic homology level}
\end{statement}
Let $C$ be the connected cyclic cover of the five-punctured sphere with
deck group $C_{12}$ and local monodromies
\[
 (9,9,10,4,4)\in(\mathbb Z/12)^5.
\]
Fill all punctures upstairs, and let $T$ denote the positive deck
transformation on the resulting integral homology. The compact curve has
genus $11$. Its degree-four and degree-three quotient covers have branch
data $(1,1,2,0,0)$ and $(0,0,1,1,1)$, respectively, and each has genus one.
Thus this single cover is a natural geometric attempt to retain the
quartic and cubic elliptic structures together.

Use the five pure braid elements
\[
 X_1=\sigma_1^2,\quad X_2=\sigma_2^2,\quad
 X_3=\sigma_3^2,\quad X_4=(\sigma_1\sigma_2\sigma_1)^2,
 \quad X_5=(\sigma_3\sigma_2\sigma_3)^2,
\]
and descend their homology action from $PB_4$ to the pure five-point
mapping class group by the linking-number normalization specified below.
For $p=2$ and for $p=3$, the closed subgroup generated by any adjacent
pair $(X_j,X_{j+1})$, with indices modulo five, has pro-$p$ second
derived subgroup in its action on $H_1(C,\mathbb Z_p)$.

Consequently the supplied actual gentle element $k$ of
\textup{FP.36}, with
\[
 f_k\in R_{\rm nil}(\widehat F(x,y)'')\qquad\hbox{and}\qquad
 \Delta(k)=1,
\]
has all five evaluations $f_k(X_j,X_{j+1})$ equal to the identity on
both $H_1(C,\mathbb Z_2)$ and $H_1(C,\mathbb Z_3)$, hence on every
finite quotient of these lattices. This assertion uses the supplied
existence statement for $k$; it does not reconstruct that lifting input.
It proves no claim about other primes or about nonabelian covering
invariants.

\begin{proof}
The sum of the branch weights is zero modulo $12$, and the weights
generate $\mathbb Z/12$, proving that the cover is well-defined and
connected. A puncture of weight $w$ has $\gcd(12,w)$ preimages, each
with ramification index $12/\gcd(12,w)$. The five numbers of preimages
are $3,3,2,4,4$. Removing the five branch values gives Euler
characteristic $12(2-5)=-36$; adding the $16$ points upstairs gives
$\chi(C)=-20$, so its genus is $11$. The same count at degrees four
and three gives Euler characteristic zero for both stated quotient
covers.

Here is a complete finite cellular model, including the marking used
by the certificate. Remove the fifth point and represent the other four
inertias by the free generators $x_1,x_2,x_3,x_4$. On the degree-twelve
cover of the resulting four-petal rose put
\[
 C_0=\mathbb Z^{12},\qquad C_1=\mathbb Z^{48},\qquad
 \partial e_{j,g}=v_{g+w_j}-v_g,
 \quad(w_1,w_2,w_3,w_4)=(9,9,10,4).
\]
Here $g\in\mathbb Z/12$, and $T(e_{j,g})=e_{j,g+1}$.
The graph is connected, so $\ker\partial$ has rank $37$.
The compact surface is obtained by filling all lifted peripheral
circles: for the first four punctures these are the translates of
$x_j^{12/\gcd(12,w_j)}$, and for the fifth they are the translates of
$(x_1x_2x_3x_4)^3$. Choose one translate for every coset of
$\langle w_j\rangle$ or $\langle8\rangle$, respectively.
There are sixteen such circles, with one relation, so their image has
rank fifteen and the quotient homology has rank $22$.
Equivalently, filling the boundary circles of a connected orientable
surface gives precisely their sole total-boundary relation. This also
shows that the quotient lattice is torsion-free.

The positive Artin automorphism is
\[
 \sigma_i(x_i)=x_ix_{i+1}x_i^{-1},\qquad
 \sigma_i(x_{i+1})=x_i,
\]
and fixes the other generators. The usual braid relations follow by
substitution in this free group; the certificate also checks all three
relations relevant to $B_4$ as literal reduced words. Each pure braid
preserves the weight map, acts on the displayed graph chains, preserves
the peripheral submodule, and commutes with $T$.

Write $P=x_1x_2x_3x_4$. The full twist
$z_4=(\sigma_1\sigma_2\sigma_3)^4$ acts by conjugation by $P$.
On the homology of the kernel of the weight map, this is the deck
action $T^8$. This identity follows because $P$ has weight $8$ and
changing a lift by an element of the kernel changes the action only
by an inner automorphism, which is trivial on its abelianization.
Let $\ell_{12}:PB_4\to\mathbb Z$ be the first-two-strands linking
number. Since $\ell_{12}(z_4)=1$, the corrected representation
\[
 \gamma\longmapsto T^{-8\ell_{12}(\gamma)}q(\gamma)
\]
is trivial on $z_4$ and descends to $PB_4/\langle z_4\rangle$.
For the five displayed elements the linking numbers are
$(1,0,0,1,0)$, so the correction multiplies precisely $X_1$ and
$X_4$ by $T^4$. Because the deck action is central, this correction
does not change the evaluation of any word $f\in\widehat F'$ at a
pair of elements.

Reduce the homology modulo $p$. The semisimple submodule, or socle,
for the deck action is
\[
 S_p=\bigoplus_h\ker h(T),
\]
where $h$ runs through the distinct monic irreducible factors of
$T^{12}-1$ in $\mathbb F_p[T]$. For $p=2$ these are
$T-1,T^2+T+1$; for $p=3$ they are $T-1,T+1,T^2+1$.
The chain calculation just described gives the following table. Each
entry is the sequence of orders of the successive derived groups of
the indicated adjacent-pair image on the indicated socle summand;
the last group is trivial.
\[
\begin{array}{c|c|c|ccccc}
 p&h&\dim_{\mathbb F_p}\ker h(T)&12&23&34&45&51\\\hline
2&T-1&1&1&1&1&1&1\\
2&T^2+T+1&6&(9,1)&(36,4,1)&(36,4,1)&(36,4,1)&(9,1)\\
3&T-1&1&1&1&1&1&1\\
3&T+1&3&(2,1)&(2,1)&1&(2,1)&(2,1)\\
3&T^2+1&6&(36,9,1)&(72,9,1)&(2,1)&(4,1)&(72,9,1)
\end{array}
\]
In particular every adjacent-pair image on $S_p$ is metabelian.

For reproducibility, \texttt{cyclic12\_cover.py} builds the boundary
matrix and all sixteen peripheral columns over the integers, computes
the compact quotient separately over $\mathbb F_2$ and $\mathbb F_3$,
and derives each pure braid matrix by tracing the literal Artin words
through the covering graph. It checks that every matrix preserves the
relation subspace, that the compact dimension is $22$, that $T^{12}=1$,
that the pure matrices commute with $T$, and that the full twist acts
as $T^8$. It then restricts to every irreducible deck socle summand.
For each pair it closes the orbit of a basis under the pair's matrices;
the resulting permutation action is faithful, since a linear map
fixing that orbit fixes the basis. Exact permutation derived-group
calculations produce the table. The outputs are
\texttt{cyclic12\_cover.json} and the two \texttt{.npz} matrix records.
Thus the computation starts with the specified covering and free-word
action, rather than assuming guessed braid matrices or a quotient
identification.

We next justify why this finite table controls all $p$-adic depths.
For a finite-dimensional $\mathbb F_p[T]$-module $V$, the kernel of
\[
 \operatorname{Aut}_{\mathbb F_p[T]}(V)
 \longrightarrow\operatorname{Aut}_{\mathbb F_p[T]}(\operatorname{soc}V)
\]
is a $p$-group. Indeed, if the kernel had an element of prime order
$q\ne p$, extend scalars to an algebraic closure. That element is
semisimple and, if nonidentity, has a nonzero eigenspace with eigenvalue
different from one. The eigenspace is $T$-stable because the two actions
commute. A nonzero $T$-stable space contains a $T$-eigenvector, which
lies in the scalar extension of the socle. This contradicts trivial
action on the socle. Cauchy's theorem now proves the claim.

Every adjacent-pair second derived group acts trivially on $S_p$ by
the table, so its image on $H_1(C,\mathbb F_p)$ is a $p$-group.
The kernel of reduction from $\operatorname{GL}_{22}(\mathbb Z_p)$
to $\operatorname{GL}_{22}(\mathbb F_p)$ is pro-$p$, as is seen from
the successive congruence kernels. Therefore the closed second
derived group of each $p$-adic pair image is pro-$p$.

Finally the restriction of each evaluation homomorphism
$\widehat F\to\operatorname{GL}_{22}(\mathbb Z_p)$ to
$\widehat F''$ has pro-$p$ image. Every finite continuous quotient
of that image is a finite $p$-group and hence nilpotent. It follows
directly from the definition of $R_{\rm nil}(\widehat F'')$ as the
intersection of kernels of finite nilpotent quotients that the
restriction kills $R_{\rm nil}(\widehat F'')$. Apply this to the
supplied $f_k$ for each of the five pairs and for $p=2,3$.
\end{proof}

\paragraph{Scope and source locators.}
The five pure braid elements and their pentagon order are those in
\texttt{DELTA\_ON\_GT/agent\_five/five\_point\_results.tex},
\textup{FP.30}. The sole actual-element existence input is
\textup{FP.36} in the same file, equivalently the strengthened
relative lifting witness used in \textup{A.94} and \textup{A.109}.
The cyclic covering, its chain action, the socle calculation, and the
all-depth implication above are new calculations in this run.
